\chapter{Sai để sáng}
\markboth{Sai để sáng}{Wrong to Become Clearer}

\section{Một học sinh giải sai rồi mới hiểu bài thật.}
\begin{truyen}{Một học sinh giải sai rồi mới hiểu bài thật.}{A mistake can open the lesson.}
\chuhoa{E}{m từng nghĩ sai là điều xấu hổ. Nhưng chính bài làm sai được cô sửa kỹ nhất lại là lúc em hiểu sâu nhất. Có những cái sai không nên che đi, vì nó là cửa mở vào chỗ mình chưa biết.}
\textit{(The student once thought being wrong was shameful. Yet the carefully corrected wrong answer became the moment of deepest understanding. Some mistakes should not be hidden because they open the door to what we do not know.)}
\end{truyen}
\begin{baihoc}
Sai trong lúc học còn tốt hơn giấu dốt rồi mang cái sai đi rất xa.
\textit{(Being wrong while learning is better than hiding ignorance and carrying that mistake much further.)}
\end{baihoc}
\begin{ghinhoanh}\textbf{Short English to remember:}

A mistake can open the lesson.
\end{ghinhoanh}
\begin{tuhocnhanh}
1. Em có đang ngại sai đến mức không dám hỏi không? \textit{(Are you so afraid of being wrong that you no longer ask questions?)}

2. Lỗi gần đây dạy em điều gì? \textit{(What did your recent mistake teach you?)}
\end{tuhocnhanh}
\ngancach

\section{Một người chọn sai bạn rồi mới biết quý người tử tế.}
\begin{truyen}{Một người chọn sai bạn rồi mới biết quý người tử tế.}{Wrong company sharpens discernment.}
\chuhoa{S}{au lần bị lợi dụng và bỏ rơi, cô mới phân biệt rõ sự ồn ào và sự chân thành. Trải nghiệm ấy đau, nhưng nó làm mắt cô sáng hơn khi chọn người đi cùng.}
\textit{(After being used and abandoned, she finally learned to distinguish noise from sincerity. The experience hurt, but it sharpened her eyes when choosing companions.)}
\end{truyen}
\begin{baihoc}
Có những cái sai trong chọn người làm mình tỉnh hơn rất nhiều về lòng người.
\textit{(Some mistakes in choosing people make us much more awake to human nature.)}
\end{baihoc}
\begin{ghinhoanh}\textbf{Short English to remember:}

Wrong company sharpens discernment.
\end{ghinhoanh}
\begin{tuhocnhanh}
1. Em học được gì từ người từng làm em thất vọng? \textit{(What did you learn from someone who once disappointed you?)}

2. Từ nay em sẽ chọn người dựa trên điều gì? \textit{(From now on, what will guide your choice of people?)}
\end{tuhocnhanh}
\ngancach

\section{Một người quyết sai ngành rồi mới thấy rõ mình hợp gì.}
\begin{truyen}{Một người quyết sai ngành rồi mới thấy rõ mình hợp gì.}{A wrong path can clarify the right one.}
\chuhoa{A}{nh học nhiều năm mới nhận ra mình chỉ đang đi theo mong muốn của người khác. Khoảng thời gian đó không mất hoàn toàn; nó giúp anh biết rất rõ mình không hợp điều gì, và vì thế thấy điều hợp với mình rõ hơn.}
\textit{(He studied for years before realizing he was following someone else's wishes. That time was not entirely wasted; it taught him clearly what did not fit, and therefore made the fitting path easier to see.)}
\end{truyen}
\begin{baihoc}
Đi nhầm đường một đoạn chưa phải là vô nghĩa, nếu ta đủ tỉnh để nhìn ra và đổi hướng.
\textit{(Walking the wrong road for a while is not meaningless if we are awake enough to notice and change direction.)}
\end{baihoc}
\begin{ghinhoanh}\textbf{Short English to remember:}

A wrong path can clarify the right one.
\end{ghinhoanh}
\begin{tuhocnhanh}
1. Có điều gì em đang làm chỉ vì người khác muốn? \textit{(What are you doing mainly because others want it?)}

2. Em cần dấu hiệu nào để biết đã đến lúc đổi hướng? \textit{(What sign tells you it may be time to change direction?)}
\end{tuhocnhanh}
\ngancach

\section{Một người nói sai một câu rồi biết giữ miệng hơn.}
\begin{truyen}{Một người nói sai một câu rồi biết giữ miệng hơn.}{One careless word can teach restraint.}
\chuhoa{C}{hỉ một câu bốc đồng trong lúc nóng giận đã làm người khác tổn thương sâu. Sau đó anh hiểu lời nói không phải thứ có thể lấy lại nguyên vẹn. Từ một lần lỡ miệng, anh học được sự chậm lại quý giá.}
\textit{(One impulsive sentence spoken in anger deeply hurt another person. After that, he understood that words cannot be taken back intact. From one careless moment, he learned the value of slowing down.)}
\end{truyen}
\begin{baihoc}
Một lỗi lời nói nếu được nhớ kỹ có thể cứu rất nhiều lần nói sai về sau.
\textit{(A remembered mistake in speech can save many future conversations.)}
\end{baihoc}
\begin{ghinhoanh}\textbf{Short English to remember:}

One careless word can teach restraint.
\end{ghinhoanh}
\begin{tuhocnhanh}
1. Em có đang mang tiếc nuối nào vì một câu đã nói? \textit{(Are you carrying regret over something you said?)}

2. Từ nay em muốn chậm lại ở tình huống nào? \textit{(In what situations do you want to slow down from now on?)}
\end{tuhocnhanh}
\ngancach

\section{Một người tin sai rồi mới học cách kiểm chứng.}
\begin{truyen}{Một người tin sai rồi mới học cách kiểm chứng.}{Being fooled can teach caution.}
\chuhoa{V}{ì tin nhanh, anh mất tiền và mất cả lòng tin vào người khác một thời gian. Nhưng từ đó anh học được cách kiểm tra, hỏi thêm, chậm quyết hơn. Sai một lần làm anh bớt cả tin, mà không biến anh thành người cay độc.}
\textit{(Because he trusted too quickly, he lost money and trust in others for a while. But from then on he learned to verify, ask more, and decide slower. One mistake made him less gullible without turning him bitter.)}
\end{truyen}
\begin{baihoc}
Sai khi tin người không đáng làm ta đóng tim hẳn; nó nên dạy ta sáng hơn trong niềm tin.
\textit{(Being wrong about someone should not shut the heart completely; it should make trust wiser.)}
\end{baihoc}
\begin{ghinhoanh}\textbf{Short English to remember:}

Being fooled can teach caution.
\end{ghinhoanh}
\begin{tuhocnhanh}
1. Em đang tin ai hoặc tin điều gì quá nhanh không? \textit{(Are you trusting someone or something too quickly?)}

2. Em cần kiểm chứng điều gì trước khi quyết tiếp? \textit{(What do you need to verify before moving on?)}
\end{tuhocnhanh}
\ngancach